\section{References}

% Long archival URLs cannot always justify cleanly; the reference lists are
% therefore set ragged right within this section only.
\begingroup
\raggedright

All web sources were accessed and checked on 2026-07-25. Ancient-work
citations in the text use conventional titles and loci with the
translations listed here.

\subsection*{Scripture}

\begin{itemize}
\item The Holy Bible, Douay--Rheims version, Challoner revision, as
transcribed at drbo.org (public domain). Chapters checked verbatim for
this study: Mt 4, 10, 14, 16, 17, 26; Mk 1, 8, 14; Lk 5, 22, 24; Jn 1, 6,
13, 18, 20, 21; Acts 1--5, 8, 10, 12, 15; 1~Cor 1, 9, 15; Gal 1--2; 1~Pt
1, 5; 2~Pt 1, 3 (chapter pages under
\url{https://drbo.org/chapter/47016.htm} and the same numbering pattern,
e.g.\ \texttt{50021} for John 21, \texttt{55002} for Galatians 2).
\end{itemize}

\subsection*{Ancient witnesses}

\begin{itemize}
\item Clement of Rome (attributed), \work{First Epistle to the
Corinthians} 5--6, in the translation presented at New Advent and there
attributed to John Keith, \textit{Ante-Nicene Fathers}, vol.~9 (Buffalo,
1896--1897) (\url{https://www.newadvent.org/fathers/1010.htm}).
Conventional dating c.~95--96 as reported in the 1912 Catholic
Encyclopedia, ``Pope St.~Clement I''
(\url{https://www.newadvent.org/cathen/04012c.htm}).
\item Ignatius of Antioch, \work{Epistle to the Romans} 4, middle
(``shorter'') recension, trans.\ A.~Roberts and J.~Donaldson, ANF vol.~1
(Buffalo, 1885), at New Advent
(\url{https://www.newadvent.org/fathers/0107.htm}).
\item Irenaeus, \work{Against Heresies} 3.1.1 and 3.3.2--3, trans.\
A.~Roberts and W.~Rambaut, ANF vol.~1, at New Advent
(\url{https://www.newadvent.org/fathers/0103301.htm};
\url{https://www.newadvent.org/fathers/0103303.htm}). A catalog-level
source-library record of ANF vol.~1 exists in the repository and is bound
in \texttt{research/source-bindings.toml}.
\item Tertullian, \work{The Prescription Against Heretics} 36 (trans.\
P.~Holmes) and \work{Scorpiace} 15 (trans.\ S.~Thelwall), ANF vol.~3, at
New Advent (\url{https://www.newadvent.org/fathers/0311.htm};
\url{https://www.newadvent.org/fathers/0318.htm}).
\item Eusebius of Caesarea, \work{Church History} 2.15, 2.25, 3.1, 3.3,
3.39, trans.\ A.~C. McGiffert, NPNF series~2, vol.~1 (New York, 1890), at
New Advent (\url{https://www.newadvent.org/fathers/250102.htm};
\url{https://www.newadvent.org/fathers/250103.htm}) and CCEL
(\url{https://ccel.org/ccel/schaff/npnf201.iii.viii.xxxix.html});
preserving Papias, Dionysius of Corinth, Gaius, and Origen at the loci
cited in the text.
\item Jerome, \work{On Illustrious Men} 1, trans.\ E.~C. Richardson, NPNF
series~2, vol.~3, at New Advent
(\url{https://www.newadvent.org/fathers/2708.htm}).
\item Tacitus, \work{Annals} 15.44, trans.\ A.~J. Church and W.~J.
Brodribb (1876), at Wikisource
(\url{https://en.wikisource.org/wiki/The_Annals_(Tacitus)/Book_15}).
\end{itemize}

\subsection*{Apocrypha and later tradition}

\begin{itemize}
\item \work{The Acts of Peter} (apocryphal; Greek, Asia Minor, ``not
later than A.D.\ 200'' per the translator's headnote), chs.\ 35--40 for
the \textit{Quo vadis} and martyrdom, in M.~R. James, \work{The
Apocryphal New Testament} (Oxford: Clarendon, 1924), as transcribed at
Early Christian Writings
(\url{https://www.earlychristianwritings.com/text/actspeter.html}).
\item ``St.~Peter, Prince of the Apostles,'' \textit{The Catholic
Encyclopedia}, vol.~11 (New York, 1911), at New Advent
(\url{https://www.newadvent.org/cathen/11744a.htm}); source for the 258
Catacombs notice, the feast history, and the judgment that the
25-year-episcopate notice's origin is untraceable.
\item ``Chair of Peter,'' \textit{The Catholic Encyclopedia}, vol.~3 (New
York, 1908), at New Advent
(\url{https://www.newadvent.org/cathen/03551e.htm}); the chair feasts and
the relic-chair description.
\end{itemize}

\subsection*{Archaeology of the tomb}

\begin{itemize}
\item E.~R. Smothers, S.J., ``The Bones of St.~Peter,'' \textit{Theological
Studies} 27 (1966) 79--88, read in full in the publisher's open PDF
(\url{https://theologicalstudies.net/wp-content/uploads/2022/08/27.1.4.pdf});
the principal checked scholarly control for the excavation report
(Apollonj Ghetti--Ferrua--Josi--Kirschbaum, \work{Esplorazioni}, 1951),
von Gerkan's and Prandi's stratigraphy, Guarducci's graffiti volumes and
\work{Le reliquie di Pietro} (1965), the Kaas--Segoni removal, and
Correnti's anthropology---none of which was inspected directly.
\item Paul VI, General Audience of 26 June 1968, Italian text, Vatican
website
(\url{https://www.vatican.va/content/paul-vi/it/audiences/1968/documents/hf_p-vi_aud_19680626.html});
source of the quoted announcement and its qualifications, including the
retrospective description of Pius XII's 1950 answer on the bones as
suspensive.
\item ``A striking discovery,'' \textit{L'Osservatore Romano}, English
edition, June 2024
(\url{https://www.osservatoreromano.va/en/news/2024-06/ing-026/a-striking-discovery.html});
checked for the official narrative of the excavations, Pius XII's 1950
radio-message announcement, and the aedicula dating.
\item ``The Necropolis,'' official presentation, St.~Peter's Basilica
website (\url{https://www.basilicasanpietro.va/en/products/the-necropolis});
the current official shrine claim, cited as such.
\end{itemize}

\subsection*{Ecclesial reception}

\begin{itemize}
\item First Vatican Council, dogmatic constitution \work{Pastor aeternus}
(18 July 1870), official Latin text, Vatican website
(\url{https://www.vatican.va/archive/hist_councils/i-vatican-council/documents/vat-i_const_18700718_pastor-aeternus_la.html});
verified byte-identical to the repository source-library artifact
(SHA-256 \texttt{032a5f65\ldots}) on 2026-07-25 and bound in
\texttt{research/source-bindings.toml}.
\item General Roman Calendar (\textit{Calendarium Romanum Generale},
Missale Romanum, editio typica tertia, 2002): identified through the
repository's source-library record; feast statements in the text also
rest on the Catholic Encyclopedia articles above and carry the 2026-07-25
currentness cutoff.
\end{itemize}

\endgroup
